% Generated draft: requires mathematical and source review.
\documentclass[11pt]{article}
\usepackage[margin=25mm]{geometry}
\usepackage{fontspec}
\setmainfont{Latin Modern Roman}
\usepackage{amsmath,amssymb,mathtools,unicode-math}
\setmathfont{Latin Modern Math}
\usepackage{enumitem,needspace,xurl,hyperref}
\setlength{\parindent}{0pt}
\setlength{\parskip}{5pt}
\setlength{\emergencystretch}{3em}
\newcommand{\N}{\mathbb{N}}
\newcommand{\Z}{\mathbb{Z}}
\newcommand{\Q}{\mathbb{Q}}
\newcommand{\R}{\mathbb{R}}
\newcommand{\C}{\mathbb{C}}
\newcommand{\F}{\mathbb{F}}
\newcommand{\A}{\mathbb{A}}
\newcommand{\PP}{\mathbb{P}}
\newcommand{\E}{\mathbb{E}}
\newcommand{\eps}{\varepsilon}
\newcommand{\dd}{\,\mathrm{d}}
\begin{document}
\section*{UnsolvedMath ID: 30000919}
\textbf{Source problem code:} catalogue entry
\subsection{Definitions and mathematical statement}
In the usual Banach-space range \(1\leq p\leq\infty\), let \(L^p[0,1]\) be the Lebesgue space on the unit interval, and define the Volterra operator by
\[
(Vf)(x)=\int_0^x f(s)\,\dd s.
\]
Let \(\mathcal B(L^p[0,1])\) be the bounded linear operators on that space. The commutant is
\[
[V]_p=\{T\in\mathcal B(L^p[0,1]):TV=VT\}.
\]
For each \(p\in(1,\infty)\setminus\{2\}\), give an explicit characterization of \([V]_p\), with necessary and sufficient conditions for a bounded operator to belong to it. Merely repeating the commutation identity is the definition, rather than the requested characterization.

The source singles out \(p=1,2,\infty\) as known cases and describes the other values as unresolved. This is retained as a supplied status statement, not asserted as independently verified current knowledge. The source does not explicitly state the overall range of \(p\) or the scalar field. The displayed target uses the conventional \(L^p\) Banach-space interpretation; it does not silently extend the problem to quasi-Banach exponents \(0<p<1\), where the given integral is not defined on every \(L^p\) function. Real versus complex scalars and any desired form of operator representation should be confirmed from the report.
\subsection{Short English statement}
Characterize all bounded operators commuting with the Volterra integration operator on \(L^p[0,1]\) for \(1<p<\infty\), \(p\ne2\), beyond the endpoint and Hilbert-space cases reported known.
\subsection{Sources}
\begin{itemize}
\item Local UnsolvedMath registry, permanent ID \(30000919\), source code catalogue entry, ``Commutants of Volterra Operators on \(L^p\) Spaces.'' No individual proposer is supplied.
\item \emph{Mini-Workshop: Complex Approximation and Universality} (2008), Oberwolfach Reports, DOI 10.4171/owr/2008/06, pp. 297--346: \url{https://doi.org/10.4171/owr/2008/06}.
\item Supplied reference, ``He--Zhu'': \url{https://doi.org/10.1016/j.jmaa.2021.125855}.
\item Supplied literature triage, checked 2026-08-21: status ``partially solved.'' It reports the commutant solved for \(p=1,2,\infty\), the remaining \(L^p\) question unresolved, and the Volterra operator not strictly cyclic on \(L^2\). The strict-cyclicity assertion is background rather than the requested characterization. This dated assessment derives from supplied OWR curation; time-sensitive status and cited claims require independent verification.
\end{itemize}
\end{document}
